On se place dans un repère orthonormé \Oijk de l'espace.

On considère les points A(1~;~0~;~3), B$( -2~;~1~;~2)$ et C(0~;~3~;~2).

\medskip

\begin{enumerate}
\item 
	\begin{enumerate}
		\item Montrer que les points A, B et C ne sont pas alignés.
		\item Soit $\vect{n}$ le vecteur de coordonnées $\begin{pmatrix}-1\\1\\4\end{pmatrix}$Vérifier que le vecteur $\vect{n}$ est orthogonal au plan (ABC).
		\item En déduire que le plan (ABC) admet pour équation cartésienne 

$-x + y + 4z - 11 = 0$.
	\end{enumerate}
\end{enumerate}

On considère le plan $\mathcal{P}$ d'équation cartésienne $3x - 3y + 2z - 9 = 0$ et le plan $\mathcal{P}'$ d'équation cartésienne $x - y- z + 2 = 0$.

\begin{enumerate}[resume]
\item 
	\begin{enumerate}
		\item Démontrer que les plans $\mathcal{P}$ et $\mathcal{P}'$ sont sécants. On note $(d)$ leur droite d'intersection.
		\item Déterminer si les plans $\mathcal{P}$ et $\mathcal{P}'$ sont perpendiculaires.
	\end{enumerate}
\item Montrer que la droite $(d)$ est dirigée par le vecteur $\vect{u}\begin{pmatrix}1\\1\\0\end{pmatrix}$.
\item Montrer que le point M(2~;~1~;~3) appartient aux plans $\mathcal{P}$ et $\mathcal{P}'$. En déduire une représentation paramétrique de la droite $(d)$.
\item Montrer que la droite $(d)$ est aussi incluse dans le plan (ABC).

Que peut-on dire des trois plans (ABC), $\mathcal{P}$ et $\mathcal{P}'$ ?
\end{enumerate}
